%lav margins mindre (bedre til normale afleveringer)
\setlrmarginsandblock{2.5cm}{*}{1} \setulmarginsandblock{3cm}{3cm}{*} \checkandfixthelayout

\usepackage[utf8]{inputenc}  % Korrekt håndtering af æ, ø og å
\usepackage[T1]{fontenc}  % Korrekt håndtering af æ, ø og å
\usepackage{microtype}  % Typografi! Giver bl.a. pænere orddeling
\usepackage[english]{babel}  % Danske betegnelser og orddeling

\usepackage[notext]{kpfonts}
\usepackage{amsmath, amssymb, bm, mathtools}  % Giver adgang til matematikting

\usepackage{graphicx}  % Gør det muligt at indsætte billeder
\usepackage{float} % gør det muligt at tvinge en figurs placering med [H]
\usepackage{listings} %Gør det muligt at indsætte pæn code
\usepackage[small,bf,hang]{caption}  %Gør caption bold
\usepackage{subcaption} % Gør det muligt at lave delfigurer
\usepackage{dsfont} % giver adgang til\mathds{#} for at lave bb numre
\usepackage[margin,draft]{fixme} % allows to add fixme notes using \fxnote or \fxwarn
% TikZ
\usepackage{tikz} % Bruges til at lave flotte figurer
\usetikzlibrary{calc} % bruges til at lave nem placering i tikz
\usepackage{pgfplots} % bruges til plots i TikZ
\pgfplotsset{compat=1.18}

\usepackage{minted}
\usemintedstyle{friendly}

%%%%%%%%%%%%%%%%%%%%%%%%%%%%%%%%
% Begin enviorment setup			 %
%%%%%%%%%%%%%%%%%%%%%%%%%%%%%%%%
\usepackage{amsthm}

\usepackage{thmtools}
\usepackage[many]{tcolorbox}
\usepackage{etoolbox}
\makeatletter

\def\renewtheorem#1{%
  \expandafter\let\csname#1\endcsname\relax
  \expandafter\let\csname c@#1\endcsname\relax
  \gdef\renewtheorem@envname{#1}
  \renewtheorem@secpar
}
\def\renewtheorem@secpar{\@ifnextchar[{\renewtheorem@numberedlike}{\renewtheorem@nonumberedlike}}
\def\renewtheorem@numberedlike[#1]#2{\newtheorem{\renewtheorem@envname}[#1]{#2}}
\def\renewtheorem@nonumberedlike#1{
  \def\renewtheorem@caption{#1}
  \edef\renewtheorem@nowithin{\noexpand\newtheorem{\renewtheorem@envname}{\renewtheorem@caption}}
  \renewtheorem@thirdpar
}
\def\renewtheorem@thirdpar{\@ifnextchar[{\renewtheorem@within}{\renewtheorem@nowithin}}
\def\renewtheorem@within[#1]{\renewtheorem@nowithin[#1]}

\makeatother

\usepackage[framemethod=TikZ]{mdframed}
\mdfsetup{skipabove=1em,skipbelow=0em, innertopmargin=12pt, innerbottommargin=8pt}

\tcbuselibrary{skins}

% Color definitions


% Dark orange and Dark Red rgb
\definecolor{theorembordercolor}{RGB}{151, 63, 5}
\definecolor{theorembackgroundcolor}{RGB}{248, 241, 234}

\definecolor{examplebordercolor}{RGB}{0, 110, 184}
\definecolor{examplebackgroundcolor}{RGB}{240, 244, 250}

\definecolor{definitionbordercolor}{RGB}{0, 150, 85}
\definecolor{definitionbackgroundcolor}{RGB}{239, 247, 243}

\definecolor{propertybordercolor}{RGB}{128, 0, 128}
\definecolor{propertybackgroundcolor}{RGB}{255, 240, 255}

\definecolor{formulabordercolor}{RGB}{0, 0, 0}
\definecolor{formulabackgroundcolor}{RGB}{230, 229, 245}
\definecolor{proofcolor}{RGB}{0,0,0}

% macro that will hold the *name* of the colour to use for the next proof
\newcommand{\lastproofcolor}{proofcolor} % initial value: the name "proofcolor"


%%% THEOREM STYLES SETUP %%%
\newtheoremstyle{theorem} % name of the style to be used
{0pt} % measure of space to leave above the theorem. E.g.: 3pt
{0pt} % measure of space to leave below the theorem. E.g.: 3pt
{\normalfont} % name of font to use in the body of the theorem
{0pt} % measure of space to indent
{\bfseries} % name of head font
{\\} % punctuation between head and body
{0pt} % space after theorem head; " " = normal interword space
{{\color{theorembordercolor}\thmname{#1}~\thmnumber{\textup{#2}}~\textup{#3}}} % Header display 
\newtheoremstyle{theorem*} % name of the style to be used
{0pt} % measure of space to leave above the theorem. E.g.: 3pt
{0pt} % measure of space to leave below the theorem. E.g.: 3pt
{\normalfont} % name of font to use in the body of the theorem
{0pt} % measure of space to indent
{\bfseries} % name of head font
{\\} % punctuation between head and body
{0pt} % space after theorem head; " " = normal interword space
{{\color{theorembordercolor}\thmname{#1}~\textup{#3}}} % Header display 


\newtheoremstyle{definition} % name of the style to be used
{0pt} % measure of space to leave above the theorem. E.g.: 3pt
{0pt} % measure of space to leave below the theorem. E.g.: 3pt
{\normalfont} % name of font to use in the body of the theorem
{0pt} % measure of space to indent
{\bfseries} % name of head font
{\\} % punctuation between head and body
{0pt} % space after theorem head; " " = normal interword space
{{\color{definitionbordercolor}\thmname{#1}~\thmnumber{\textup{#2}}~\textup{#3}}} % Header display 
\newtheoremstyle{definition*} % name of the style to be used
{0pt} % measure of space to leave above the theorem. E.g.: 3pt
{0pt} % measure of space to leave below the theorem. E.g.: 3pt
{\normalfont} % name of font to use in the body of the theorem
{0pt} % measure of space to indent
{\bfseries} % name of head font
{\\} % punctuation between head and body
{0pt} % space after theorem head; " " = normal interword space
{{\color{definitionbordercolor}\thmname{#1}~\textup{#3}}} % Header display 

\newtheoremstyle{example} % name of the style to be used
{0pt} % measure of space to leave above the theorem. E.g.: 3pt
{0pt} % measure of space to leave below the theorem. E.g.: 3pt
{\normalfont} % name of font to use in the body of the theorem
{0pt} % measure of space to indent
{\bfseries} % name of head font
{\\} % punctuation between head and body
{0pt} % space after theorem head; " " = normal interword space
{{\color{examplebordercolor}\thmname{#1}~\thmnumber{\textup{#2}}~\textup{#3}}} % Header display 
\newtheoremstyle{example*} % name of the style to be used
{0pt} % measure of space to leave above the theorem. E.g.: 3pt
{0pt} % measure of space to leave below the theorem. E.g.: 3pt
{\normalfont} % name of font to use in the body of the theorem
{0pt} % measure of space to indent
{\bfseries} % name of head font
{\\} % punctuation between head and body
{0pt} % space after theorem head; " " = normal interword space
{{\color{examplebordercolor}\thmname{#1}~\textup{#3}}} % Header display 

\newtheoremstyle{formula} % name of the style to be used
{0pt} % measure of space to leave above the theorem. E.g.: 3pt
{0pt} % measure of space to leave below the theorem. E.g.: 3pt
{\normalfont} % name of font to use in the body of the theorem
{0pt} % measure of space to indent
{\bfseries} % name of head font
{\\} % punctuation between head and body
{0pt} % space after theorem head; " " = normal interword space
{{\color{formulabordercolor}\thmname{#1}~\thmnumber{\textup{#2}}~\textup{#3}}} % Header display 
\newtheoremstyle{formula*} % name of the style to be used
{0pt} % measure of space to leave above the theorem. E.g.: 3pt
{0pt} % measure of space to leave below the theorem. E.g.: 3pt
{\normalfont} % name of font to use in the body of the theorem
{0pt} % measure of space to indent
{\bfseries} % name of head font
{\\} % punctuation between head and body
{0pt} % space after theorem head; " " = normal interword space
{{\color{formulabordercolor}\thmname{#1}~\textup{#3}}} % Header display 


\newtheoremstyle{example} % name of the style to be used
{0pt} % measure of space to leave above the theorem. E.g.: 3pt
{0pt} % measure of space to leave below the theorem. E.g.: 3pt
{\normalfont} % name of font to use in the body of the theorem
{0pt} % measure of space to indent
{\bfseries} % name of head font
{\\} % punctuation between head and body
{0pt} % space after theorem head; " " = normal interword space
{{\color{examplebordercolor}\thmname{#1}~\thmnumber{\textup{#2}}~\textup{#3}}} % Header display 
\newtheoremstyle{example*} % name of the style to be used
{0pt} % measure of space to leave above the theorem. E.g.: 3pt
{0pt} % measure of space to leave below the theorem. E.g.: 3pt
{\normalfont} % name of font to use in the body of the theorem
{0pt} % measure of space to indent
{\bfseries} % name of head font
{\\} % punctuation between head and body
{0pt} % space after theorem head; " " = normal interword space
{{\color{examplebordercolor}\thmname{#1}~\textup{#3}}} % Header display 

\newtheoremstyle{property} % name of the style to be used
{0pt} % measure of space to leave above the theorem. E.g.: 3pt
{0pt} % measure of space to leave below the theorem. E.g.: 3pt
{\normalfont} % name of font to use in the body of the theorem
{0pt} % measure of space to indent
{\bfseries} % name of head font
{\\} % punctuation between head and body
{0pt} % space after theorem head; " " = normal interword space
{{\color{propertybordercolor}\thmname{#1}~\thmnumber{\textup{#2}}~\textup{#3}}} % Header display 
\newtheoremstyle{property*} % name of the style to be used
{0pt} % measure of space to leave above the theorem. E.g.: 3pt
{0pt} % measure of space to leave below the theorem. E.g.: 3pt
{\normalfont} % name of font to use in the body of the theorem
{0pt} % measure of space to indent
{\bfseries} % name of head font
{\\} % punctuation between head and body
{0pt} % space after theorem head; " " = normal interword space
{{\color{propertybordercolor}\thmname{#1}~\textup{#3}}} % Header display 


\newtheoremstyle{question} % name of the style to be used
{0pt} % measure of space to leave above the theorem. E.g.: 3pt
{0pt} % measure of space to leave below the theorem. E.g.: 3pt
{\normalfont} % name of font to use in the body of the theorem
{0pt} % measure of space to indent
{\bfseries} % name of head font
{\\} % punctuation between head and body
{0pt} % space after theorem head; " " = normal interword space
{{\color{formulabordercolor}\thmname{#1}~\textup{#3}}} % Header display 

%%%%%%%%%%%%%%%%%%%%%%%%
% Theorem Environments 				%
%%%%%%%%%%%%%%%%%%%%%%%%

\theoremstyle{theorem}
\newtheorem{theorem}{Theorem}
\newtheorem{postulate}[theorem]{Postulate}
\newtheorem{conjecture}[theorem]{Conjecture}
\newtheorem{corollary}[theorem]{Corollary}
\newtheorem{lemma}[theorem]{Lemma}
\newtheorem{conclusion}[theorem]{Conclusion}
% unnumbered versions
\theoremstyle{theorem*}
\newtheorem{theorem*}{Theorem}
\newtheorem{postulate*}{Postulate}
\newtheorem{conjecture*}{Conjecture}
\newtheorem{corollary*}{Corollary}
\newtheorem{lemma*}{Lemma}
\newtheorem{conclusion*}{Conclusion}

\tcolorboxenvironment{theorem}{
  enhanced jigsaw, pad at break*=1mm, breakable,
  left=4mm, right=4mm, top=1mm, bottom=1mm,
  colback=theorembackgroundcolor, boxrule=0pt, frame hidden,
  borderline west={0.5mm}{0mm}{theorembordercolor}, arc=.5mm
}
\tcolorboxenvironment{postulate}{
  enhanced jigsaw, pad at break*=1mm, breakable,
  left=4mm, right=4mm, top=1mm, bottom=1mm,
  colback=theorembackgroundcolor, boxrule=0pt, frame hidden,
  borderline west={0.5mm}{0mm}{theorembordercolor}, arc=.5mm
}
\tcolorboxenvironment{conjecture}{
  enhanced jigsaw, pad at break*=1mm, breakable,
  left=4mm, right=4mm, top=1mm, bottom=1mm,
  colback=theorembackgroundcolor, boxrule=0pt, frame hidden,
  borderline west={0.5mm}{0mm}{theorembordercolor}, arc=.5mm
}
\tcolorboxenvironment{lemma}{
  enhanced jigsaw, pad at break*=1mm, breakable,
  left=4mm, right=4mm, top=1mm, bottom=1mm,
  colback=theorembackgroundcolor, boxrule=0pt, frame hidden,
  borderline west={0.5mm}{0mm}{theorembordercolor}, arc=.5mm
}
\tcolorboxenvironment{conclusion}{
  enhanced jigsaw, pad at break*=1mm, breakable,
  left=4mm, right=4mm, top=1mm, bottom=1mm,
  colback=theorembackgroundcolor, boxrule=0pt, frame hidden,
  borderline west={0.5mm}{0mm}{theorembordercolor}, arc=.5mm
}
\tcolorboxenvironment{corollary}{
  enhanced jigsaw, pad at break*=1mm, breakable,
  left=4mm, right=4mm, top=1mm, bottom=1mm,
  colback=theorembackgroundcolor, boxrule=0pt, frame hidden,
  borderline west={0.5mm}{0mm}{theorembordercolor}, arc=.5mm
}

\tcolorboxenvironment{theorem*}{
  enhanced jigsaw, pad at break*=1mm, breakable,
  left=4mm, right=4mm, top=1mm, bottom=1mm,
  colback=theorembackgroundcolor, boxrule=0pt, frame hidden,
  borderline west={0.5mm}{0mm}{theorembordercolor}, arc=.5mm
}
\tcolorboxenvironment{postulate*}{
  enhanced jigsaw, pad at break*=1mm, breakable,
  left=4mm, right=4mm, top=1mm, bottom=1mm,
  colback=theorembackgroundcolor, boxrule=0pt, frame hidden,
  borderline west={0.5mm}{0mm}{theorembordercolor}, arc=.5mm
}
\tcolorboxenvironment{conjecture*}{
  enhanced jigsaw, pad at break*=1mm, breakable,
  left=4mm, right=4mm, top=1mm, bottom=1mm,
  colback=theorembackgroundcolor, boxrule=0pt, frame hidden,
  borderline west={0.5mm}{0mm}{theorembordercolor}, arc=.5mm
}
\tcolorboxenvironment{lemma*}{
  enhanced jigsaw, pad at break*=1mm, breakable,
  left=4mm, right=4mm, top=1mm, bottom=1mm,
  colback=theorembackgroundcolor, boxrule=0pt, frame hidden,
  borderline west={0.5mm}{0mm}{theorembordercolor}, arc=.5mm
}
\tcolorboxenvironment{conclusion*}{
  enhanced jigsaw, pad at break*=1mm, breakable,
  left=4mm, right=4mm, top=1mm, bottom=1mm,
  colback=theorembackgroundcolor, boxrule=0pt, frame hidden,
  borderline west={0.5mm}{0mm}{theorembordercolor}, arc=.5mm
}
\tcolorboxenvironment{corollary*}{
  enhanced jigsaw, pad at break*=1mm, breakable,
  left=4mm, right=4mm, top=1mm, bottom=1mm,
  colback=theorembackgroundcolor, boxrule=0pt, frame hidden,
  borderline west={0.5mm}{0mm}{theorembordercolor}, arc=.5mm
}


%%%%%%%%%%%%%%%%%%%%%%%%%%%
% Definition Environments %
%%%%%%%%%%%%%%%%%%%%%%%%%%%

\theoremstyle{definition}
\newtheorem{definition}[theorem]{Definition}
\newtheorem{review}[theorem]{Review}

% Unnumbered version 
\theoremstyle{definition*}
\newtheorem{definition*}{Definition}
\newtheorem{review*}{Review}


\tcolorboxenvironment{definition}{
  enhanced jigsaw, pad at break*=1mm, breakable,
  left=4mm, right=4mm, top=1mm, bottom=1mm,
  colback=definitionbackgroundcolor, boxrule=0pt, frame hidden,
  borderline west={0.5mm}{0mm}{definitionbordercolor}, arc=.5mm
}
\tcolorboxenvironment{review}{
  enhanced jigsaw, pad at break*=1mm, breakable,
  left=4mm, right=4mm, top=1mm, bottom=1mm,
  colback=definitionbackgroundcolor, boxrule=0pt, frame hidden,
  borderline west={0.5mm}{0mm}{definitionbordercolor}, arc=.5mm
}

\tcolorboxenvironment{definition*}{
  enhanced jigsaw, pad at break*=1mm, breakable,
  left=4mm, right=4mm, top=1mm, bottom=1mm,
  colback=definitionbackgroundcolor, boxrule=0pt, frame hidden,
  borderline west={0.5mm}{0mm}{definitionbordercolor}, arc=.5mm
}
\tcolorboxenvironment{review*}{
  enhanced jigsaw, pad at break*=1mm, breakable,
  left=4mm, right=4mm, top=1mm, bottom=1mm,
  colback=definitionbackgroundcolor, boxrule=0pt, frame hidden,
  borderline west={0.5mm}{0mm}{definitionbordercolor}, arc=.5mm
}


%%%%%%%%%%%%%%%%%%%%%%%%
% Example Environments 			%
%%%%%%%%%%%%%%%%%%%%%%%%

\theoremstyle{example}
\newtheorem{example}[theorem]{Example}
\newtheorem{remark}[theorem]{Remark}
\newtheorem{note}[theorem]{Note}
\newtheorem{claim}[theorem]{Claim}
\newtheorem{fact}[theorem]{Fact}

\theoremstyle{example*}
\newtheorem{example*}{Example}
\newtheorem{remark*}{Remark}
\newtheorem{note*}{Note}
\newtheorem{claim*}{Claim}
\newtheorem{fact*}{Fact}

\tcolorboxenvironment{example}{
  enhanced jigsaw, pad at break*=1mm, breakable,
  left=4mm, right=4mm, top=1mm, bottom=1mm,
  colback=examplebackgroundcolor, boxrule=0pt, frame hidden,
  borderline west={0.5mm}{0mm}{examplebordercolor}, arc=.5mm
}
\tcolorboxenvironment{remark}{
  enhanced jigsaw, pad at break*=1mm, breakable,
  left=4mm, right=4mm, top=1mm, bottom=1mm,
  colback=white, boxrule=0pt, frame hidden,
  borderline west={0.5mm}{0mm}{examplebordercolor}, arc=.5mm
}
\tcolorboxenvironment{note}{
  enhanced jigsaw, pad at break*=1mm, breakable,
  left=4mm, right=4mm, top=1mm, bottom=1mm,
  colback=white, boxrule=0pt, frame hidden,
  borderline west={0.5mm}{0mm}{examplebordercolor}, arc=.5mm
}
\tcolorboxenvironment{claim}{
  enhanced jigsaw, pad at break*=1mm, breakable,
  left=4mm, right=4mm, top=1mm, bottom=1mm,
  colback=white, boxrule=0pt, frame hidden,
  borderline west={0.5mm}{0mm}{examplebordercolor}, arc=.5mm
}
\tcolorboxenvironment{fact}{
  enhanced jigsaw, pad at break*=1mm, breakable,
  left=4mm, right=4mm, top=1mm, bottom=1mm,
  colback=examplebackgroundcolor, boxrule=0pt, frame hidden,
  borderline west={0.5mm}{0mm}{examplebordercolor}, arc=.5mm
}

\tcolorboxenvironment{example*}{
  enhanced jigsaw, pad at break*=1mm, breakable,
  left=4mm, right=4mm, top=1mm, bottom=1mm,
  colback=examplebackgroundcolor, boxrule=0pt, frame hidden,
  borderline west={0.5mm}{0mm}{examplebordercolor}, arc=.5mm
}
\tcolorboxenvironment{remark*}{
  enhanced jigsaw, pad at break*=1mm, breakable,
  left=4mm, right=4mm, top=1mm, bottom=1mm,
  colback=white, boxrule=0pt, frame hidden,
  borderline west={0.5mm}{0mm}{examplebordercolor}, arc=.5mm
}
\tcolorboxenvironment{note*}{
  enhanced jigsaw, pad at break*=1mm, breakable,
  left=4mm, right=4mm, top=1mm, bottom=1mm,
  colback=white, boxrule=0pt, frame hidden,
  borderline west={0.5mm}{0mm}{examplebordercolor}, arc=.5mm
}
\tcolorboxenvironment{claim*}{
  enhanced jigsaw, pad at break*=1mm, breakable,
  left=4mm, right=4mm, top=1mm, bottom=1mm,
  colback=white, boxrule=0pt, frame hidden,
  borderline west={0.5mm}{0mm}{examplebordercolor}, arc=.5mm
}
\tcolorboxenvironment{fact*}{
  enhanced jigsaw, pad at break*=1mm, breakable,
  left=4mm, right=4mm, top=1mm, bottom=1mm,
  colback=examplebackgroundcolor, boxrule=0pt, frame hidden,
  borderline west={0.5mm}{0mm}{examplebordercolor}, arc=.5mm
}

%%%%%%%%%%%%%%%%%%%%%%%%%
% Property Environments 					%
%%%%%%%%%%%%%%%%%%%%%%%%%

\theoremstyle{property}
\newtheorem{property}[theorem]{Property}
\newtheorem{proposition}[theorem]{Proposition}
\newtheorem{result}[theorem]{Result}

\theoremstyle{property*}
\newtheorem{property*}{Property}
\newtheorem{proposition*}{Proposition}
\newtheorem{result*}{Result}


\tcolorboxenvironment{property}{
  enhanced jigsaw, pad at break*=1mm, breakable,
  left=4mm, right=4mm, top=1mm, bottom=1mm,
  colback=propertybackgroundcolor, boxrule=0pt, frame hidden,
  borderline west={0.5mm}{0mm}{propertybordercolor}, arc=.5mm
}
\tcolorboxenvironment{proposition}{
  enhanced jigsaw, pad at break*=1mm, breakable,
  left=4mm, right=4mm, top=1mm, bottom=1mm,
  colback=propertybackgroundcolor, boxrule=0pt, frame hidden,
  borderline west={0.5mm}{0mm}{propertybordercolor}, arc=.5mm
}
\tcolorboxenvironment{result}{
  enhanced jigsaw, pad at break*=1mm, breakable,
  left=4mm, right=4mm, top=1mm, bottom=1mm,
  colback=propertybackgroundcolor, boxrule=0pt, frame hidden,
  borderline west={0.5mm}{0mm}{propertybordercolor}, arc=.5mm
}

\tcolorboxenvironment{property*}{
  enhanced jigsaw, pad at break*=1mm, breakable,
  left=4mm, right=4mm, top=1mm, bottom=1mm,
  colback=propertybackgroundcolor, boxrule=0pt, frame hidden,
  borderline west={0.5mm}{0mm}{propertybordercolor}, arc=.5mm
}
\tcolorboxenvironment{proposition*}{
  enhanced jigsaw, pad at break*=1mm, breakable,
  left=4mm, right=4mm, top=1mm, bottom=1mm,
  colback=propertybackgroundcolor, boxrule=0pt, frame hidden,
  borderline west={0.5mm}{0mm}{propertybordercolor}, arc=.5mm
}
\tcolorboxenvironment{result*}{
  enhanced jigsaw, pad at break*=1mm, breakable,
  left=4mm, right=4mm, top=1mm, bottom=1mm,
  colback=propertybackgroundcolor, boxrule=0pt, frame hidden,
  borderline west={0.5mm}{0mm}{propertybordercolor}, arc=.5mm
}

%%%%%%%%%%%%%%%%%%%%%%%%%
% Formula Enviorments		%
%%%%%%%%%%%%%%%%%%%%%%%%%

\theoremstyle{formula}
\newtheorem{formula}[theorem]{Formula}

\theoremstyle{formula*}
\newtheorem{formula*}{Formula}


\tcolorboxenvironment{formula}{
  enhanced jigsaw, pad at break*=1mm, breakable,
  left=4mm, right=4mm, top=1mm, bottom=1mm,
  colback=formulabackgroundcolor, boxrule=0pt, frame hidden,
  borderline west={0.5mm}{0mm}{formulabordercolor}, arc=.5mm
}

\tcolorboxenvironment{formula*}{
  enhanced jigsaw, pad at break*=1mm, breakable,
  left=4mm, right=4mm, top=1mm, bottom=1mm,
  colback=formulabackgroundcolor, boxrule=0pt, frame hidden,
  borderline west={0.5mm}{0mm}{formulabordercolor}, arc=.5mm
}

%%%%%%%%%%%%%%%%%%%%%%%%%%%
% Question Enviorments 		%  (Unnumbered enviorments)
%%%%%%%%%%%%%%%%%%%%%%%%%%%

\theoremstyle{question}
\newtheorem*{question}{Question}
\newtheorem*{assigment}{Assignment}
\newtheorem*{exercise}{Exercise}

\tcolorboxenvironment{question}{
  enhanced jigsaw, pad at break*=1mm, breakable,
  left=4mm, right=4mm, top=1mm, bottom=1mm,
  colback=formulabackgroundcolor, boxrule=0pt, frame hidden,
  borderline west={0.5mm}{0mm}{formulabordercolor}, arc=.5mm
}
\tcolorboxenvironment{assigment}{
  enhanced jigsaw, pad at break*=1mm, breakable,
  left=4mm, right=4mm, top=1mm, bottom=1mm,
  colback=formulabackgroundcolor, boxrule=0pt, frame hidden,
  borderline west={0.5mm}{0mm}{formulabordercolor}, arc=.5mm
}
\tcolorboxenvironment{exercise}{
  enhanced jigsaw, pad at break*=1mm, breakable,
  left=4mm, right=4mm, top=1mm, bottom=1mm,
  colback=formulabackgroundcolor, boxrule=0pt, frame hidden,
  borderline west={0.5mm}{0mm}{formulabordercolor}, arc=.5mm
}

%%%%%%%%%%%%%%%%%%%%%
% Proof Enviorment	%
%%%%%%%%%%%%%%%%%%%%%

% proof box: use \lastproofcolor for border and for the body font
\newtheoremstyle{proof} % name of the style to be used
{0pt} % measure of space to leave above the theorem. E.g.: 3pt
{0pt} % measure of space to leave below the theorem. E.g.: 3pt
{\normalfont} % name of font to use in the body of the theorem
{0pt} % measure of space to indent
{\bfseries} % name of head font
{\\} % punctuation between head and body
{0pt} % space after theorem head; " " = normal interword space
{{\color{\lastproofcolor}\thmname{#1}~\textup{#3}}} % Header display

% These patches must be placed after \tcolorboxenvironment !
%% Change color of proof to match prev enviorment.
\AtEndEnvironment{theorem}{\global\def\lastproofcolor{theorembordercolor}}
\AtEndEnvironment{postulate}{\global\def\lastproofcolor{theorembordercolor}}
\AtEndEnvironment{conjecture}{\global\def\lastproofcolor{theorembordercolor}}
\AtEndEnvironment{corollary}{\global\def\lastproofcolor{theorembordercolor}}
\AtEndEnvironment{lemma}{\global\def\lastproofcolor{theorembordercolor}}
\AtEndEnvironment{conclusion}{\global\def\lastproofcolor{theorembordercolor}}

\AtEndEnvironment{theorem*}{\global\def\lastproofcolor{theorembordercolor}}
\AtEndEnvironment{postulate*}{\global\def\lastproofcolor{theorembordercolor}}
\AtEndEnvironment{conjecture*}{\global\def\lastproofcolor{theorembordercolor}}
\AtEndEnvironment{corollary*}{\global\def\lastproofcolor{theorembordercolor}}
\AtEndEnvironment{lemma*}{\global\def\lastproofcolor{theorembordercolor}}
\AtEndEnvironment{conclusion*}{\global\def\lastproofcolor{theorembordercolor}}

\AtEndEnvironment{definition}{\global\def\lastproofcolor{definitionbordercolor}}
\AtEndEnvironment{review}{\global\def\lastproofcolor{definitionbordercolor}}

\AtEndEnvironment{definition*}{\global\def\lastproofcolor{definitionbordercolor}}
\AtEndEnvironment{review*}{\global\def\lastproofcolor{definitionbordercolor}}

\AtEndEnvironment{example}{\global\def\lastproofcolor{examplebordercolor}}
\AtEndEnvironment{remark}{\global\def\lastproofcolor{examplebordercolor}}
\AtEndEnvironment{note}{\global\def\lastproofcolor{examplebordercolor}}
\AtEndEnvironment{claim}{\global\def\lastproofcolor{examplebordercolor}}
\AtEndEnvironment{fact}{\global\def\lastproofcolor{examplebordercolor}}


\AtEndEnvironment{example*}{\global\def\lastproofcolor{examplebordercolor}}
\AtEndEnvironment{remark*}{\global\def\lastproofcolor{examplebordercolor}}
\AtEndEnvironment{note*}{\global\def\lastproofcolor{examplebordercolor}}
\AtEndEnvironment{claim*}{\global\def\lastproofcolor{examplebordercolor}}
\AtEndEnvironment{fact*}{\global\def\lastproofcolor{examplebordercolor}}


\AtEndEnvironment{property}{\global\def\lastproofcolor{propertybordercolor}}
\AtEndEnvironment{proposition}{\global\def\lastproofcolor{propertybordercolor}}
\AtEndEnvironment{result}{\global\def\lastproofcolor{propertybordercolor}}

\AtEndEnvironment{property*}{\global\def\lastproofcolor{propertybordercolor}}
\AtEndEnvironment{proposition*}{\global\def\lastproofcolor{propertybordercolor}}
\AtEndEnvironment{result*}{\global\def\lastproofcolor{propertybordercolor}}

\AtEndEnvironment{formula}{\global\def\lastproofcolor{formulabordercolor}}
\AtEndEnvironment{question}{\global\def\lastproofcolor{formulabordercolor}}
\AtEndEnvironment{assigment}{\global\def\lastproofcolor{formulabordercolor}}
\AtEndEnvironment{exercise}{\global\def\lastproofcolor{formulabordercolor}}

\AtEndEnvironment{formula*}{\global\def\lastproofcolor{formulabordercolor}}
\AtEndEnvironment{question*}{\global\def\lastproofcolor{formulabordercolor}}
\AtEndEnvironment{assigment*}{\global\def\lastproofcolor{formulabordercolor}}
\AtEndEnvironment{exercise*}{\global\def\lastproofcolor{formulabordercolor}}


\renewcommand{\qedsymbol}{Q.E.D.}
\let\qedsymbolMyOriginal\qedsymbol
\renewcommand{\qedsymbol}{
  \color{\lastproofcolor}\qedsymbolMyOriginal
}


\theoremstyle{proof}
\renewtheorem{proof}{Proof}
\renewtheorem{solution}{Solution}
\renewtheorem{answer}{Answer}

\tcolorboxenvironment{proof}{
  enhanced jigsaw, pad at break*=1mm, breakable,
  left=4mm, right=4mm, top=1mm, bottom=1mm,
  colback=white, boxrule=0pt, frame hidden,
  borderline west={0.5mm}{0mm}{\lastproofcolor}, arc=.5mm
}
\tcolorboxenvironment{solution}{
  enhanced jigsaw, pad at break*=1mm, breakable,
  left=4mm, right=4mm, top=1mm, bottom=1mm,
  colback=white, boxrule=0pt, frame hidden,
  borderline west={0.5mm}{0mm}{\lastproofcolor}, arc=.5mm
}
\tcolorboxenvironment{answer}{
  enhanced jigsaw, pad at break*=1mm, breakable,
  left=4mm, right=4mm, top=1mm, bottom=1mm,
  colback=white, boxrule=0pt, frame hidden,
  borderline west={0.5mm}{0mm}{\lastproofcolor}, arc=.5mm
}

% Inline Fixme inside custom theorems to fix error.
\AtBeginEnvironment{proof}{\fxsetup{layout=inline}}
\AtEndEnvironment{proof}{\fxsetup{layout=margin}}
\AtBeginEnvironment{solution}{\fxsetup{layout=inline}}
\AtEndEnvironment{solution}{\fxsetup{layout=margin}}
\AtBeginEnvironment{answer}{\fxsetup{layout=inline}}
\AtEndEnvironment{answer}{\fxsetup{layout=margin}}

\AtBeginEnvironment{theorem}{\fxsetup{layout=inline}}
\AtEndEnvironment{theorem}{\fxsetup{layout=margin}}
\AtBeginEnvironment{postulate}{\fxsetup{layout=inline}}
\AtEndEnvironment{postulate}{\fxsetup{layout=margin}}
\AtBeginEnvironment{conjecture}{\fxsetup{layout=inline}}
\AtEndEnvironment{conjecture}{\fxsetup{layout=margin}}
\AtBeginEnvironment{corollary}{\fxsetup{layout=inline}}
\AtEndEnvironment{corollary}{\fxsetup{layout=margin}}
\AtBeginEnvironment{lemma}{\fxsetup{layout=inline}}
\AtEndEnvironment{lemma}{\fxsetup{layout=margin}}
\AtBeginEnvironment{conclusion}{\fxsetup{layout=inline}}
\AtEndEnvironment{conclusion}{\fxsetup{layout=margin}}

\AtBeginEnvironment{theorem*}{\fxsetup{layout=inline}}
\AtEndEnvironment{theorem*}{\fxsetup{layout=margin}}
\AtBeginEnvironment{postulate*}{\fxsetup{layout=inline}}
\AtEndEnvironment{postulate*}{\fxsetup{layout=margin}}
\AtBeginEnvironment{conjecture*}{\fxsetup{layout=inline}}
\AtEndEnvironment{conjecture*}{\fxsetup{layout=margin}}
\AtBeginEnvironment{corollary*}{\fxsetup{layout=inline}}
\AtEndEnvironment{corollary*}{\fxsetup{layout=margin}}
\AtBeginEnvironment{lemma*}{\fxsetup{layout=inline}}
\AtEndEnvironment{lemma*}{\fxsetup{layout=margin}}
\AtBeginEnvironment{conclusion*}{\fxsetup{layout=inline}}
\AtEndEnvironment{conclusion*}{\fxsetup{layout=margin}}

\AtBeginEnvironment{definition}{\fxsetup{layout=inline}}
\AtEndEnvironment{definition}{\fxsetup{layout=margin}}
\AtBeginEnvironment{review}{\fxsetup{layout=inline}}
\AtEndEnvironment{review}{\fxsetup{layout=margin}}

\AtBeginEnvironment{definition*}{\fxsetup{layout=inline}}
\AtEndEnvironment{definition*}{\fxsetup{layout=margin}}
\AtBeginEnvironment{review*}{\fxsetup{layout=inline}}
\AtEndEnvironment{review*}{\fxsetup{layout=margin}}

\AtBeginEnvironment{example}{\fxsetup{layout=inline}}
\AtEndEnvironment{example}{\fxsetup{layout=margin}}
\AtBeginEnvironment{remark}{\fxsetup{layout=inline}}
\AtEndEnvironment{remark}{\fxsetup{layout=margin}}
\AtBeginEnvironment{note}{\fxsetup{layout=inline}}
\AtEndEnvironment{note}{\fxsetup{layout=margin}}
\AtBeginEnvironment{claim}{\fxsetup{layout=inline}}
\AtEndEnvironment{claim}{\fxsetup{layout=margin}}
\AtBeginEnvironment{fact}{\fxsetup{layout=inline}}
\AtEndEnvironment{fact}{\fxsetup{layout=margin}}

\AtBeginEnvironment{example*}{\fxsetup{layout=inline}}
\AtEndEnvironment{example*}{\fxsetup{layout=margin}}
\AtBeginEnvironment{remark*}{\fxsetup{layout=inline}}
\AtEndEnvironment{remark*}{\fxsetup{layout=margin}}
\AtBeginEnvironment{note*}{\fxsetup{layout=inline}}
\AtEndEnvironment{note*}{\fxsetup{layout=margin}}
\AtBeginEnvironment{claim*}{\fxsetup{layout=inline}}
\AtEndEnvironment{claim*}{\fxsetup{layout=margin}}
\AtBeginEnvironment{fact*}{\fxsetup{layout=inline}}
\AtEndEnvironment{fact*}{\fxsetup{layout=margin}}

\AtBeginEnvironment{example*}{\fxsetup{layout=inline}}
\AtEndEnvironment{example*}{\fxsetup{layout=margin}}
\AtBeginEnvironment{remark*}{\fxsetup{layout=inline}}
\AtEndEnvironment{remark*}{\fxsetup{layout=margin}}
\AtBeginEnvironment{note*}{\fxsetup{layout=inline}}
\AtEndEnvironment{note*}{\fxsetup{layout=margin}}
\AtBeginEnvironment{claim*}{\fxsetup{layout=inline}}
\AtEndEnvironment{claim*}{\fxsetup{layout=margin}}
\AtBeginEnvironment{fact*}{\fxsetup{layout=inline}}
\AtEndEnvironment{fact*}{\fxsetup{layout=margin}}

\AtBeginEnvironment{property}{\fxsetup{layout=inline}}
\AtEndEnvironment{property}{\fxsetup{layout=margin}}
\AtBeginEnvironment{proposition}{\fxsetup{layout=inline}}
\AtEndEnvironment{proposition}{\fxsetup{layout=margin}}
\AtBeginEnvironment{result}{\fxsetup{layout=inline}}
\AtEndEnvironment{result}{\fxsetup{layout=margin}}

\AtBeginEnvironment{formula}{\fxsetup{layout=inline}}
\AtEndEnvironment{formula}{\fxsetup{layout=margin}}
\AtBeginEnvironment{question}{\fxsetup{layout=inline}}
\AtEndEnvironment{question}{\fxsetup{layout=margin}}
\AtBeginEnvironment{assigment}{\fxsetup{layout=inline}}
\AtEndEnvironment{assigment}{\fxsetup{layout=margin}}
\AtBeginEnvironment{exercise}{\fxsetup{layout=inline}}
\AtEndEnvironment{exercise}{\fxsetup{layout=margin}}

\AtBeginEnvironment{formula*}{\fxsetup{layout=inline}}
\AtEndEnvironment{formula*}{\fxsetup{layout=margin}}
\AtBeginEnvironment{question*}{\fxsetup{layout=inline}}
\AtEndEnvironment{question*}{\fxsetup{layout=margin}}
\AtBeginEnvironment{assigment*}{\fxsetup{layout=inline}}
\AtEndEnvironment{assigment*}{\fxsetup{layout=margin}}
\AtBeginEnvironment{exercise*}{\fxsetup{layout=inline}}
\AtEndEnvironment{exercise*}{\fxsetup{layout=margin}}

%%%%%%%%%%%%%%%%%%%%%%%%%%%%
% Referencing setup				 %
%%%%%%%%%%%%%%%%%%%%%%%%%%%%
% setup hyperref
\usepackage[hidelinks]{hyperref}
\hypersetup{linkcolor=black, filecolor=magenta, urlcolor=cyan, pdftitle={}, pdfpagemode=FullScreen, plainpages=false }


% Make numbering follow the chapter.
\numberwithin{equation}{chapter}
\numberwithin{figure}{chapter}
\numberwithin{table}{chapter}

\numberwithin{theorem}{chapter}
\numberwithin{postulate}{chapter}
\numberwithin{conjecture}{chapter}
\numberwithin{corollary}{chapter}
\numberwithin{lemma}{chapter}
\numberwithin{conclusion}{chapter}
\numberwithin{definition}{chapter}
\numberwithin{review}{chapter}
\numberwithin{example}{chapter}
\numberwithin{note}{chapter}
\numberwithin{claim}{chapter}
\numberwithin{fact}{chapter}
\numberwithin{property}{chapter}
\numberwithin{proposition}{chapter}
\numberwithin{result}{chapter}
\numberwithin{formula}{chapter}

% Setup cleverref to custom enviorments
\usepackage[noabbrev, nameinlink,]{cleveref}
\crefname{postulate}{Postulate}{Postulates}
\crefname{conjecture}{Conjecture}{Conjectures}
\crefname{corollary}{Corollary}{Corollaries}
\crefname{lemma}{Lemma}{Lemmas}
\crefname{conclusion}{Conclusion}{Conclusions}
\crefname{definition}{Definition}{Definitions}
\crefname{review}{Review}{Reviews}
\crefname{example}{Example}{Examples}
\crefname{note}{Note}{Notes}
\crefname{claim}{Claim}{Claims}
\crefname{fact}{Fact}{Facts}
\crefname{property}{Property}{Properties}
\crefname{proposition}{Proposition}{Propositions}
\crefname{formula}{Formula}{Formulas}

%%%%%%%%%%%%%%%%%%%%%%%%%
% Header & Footer setup %
%%%%%%%%%%%%%%%%%%%%%%%%%
\usepackage{fancyhdr}
\pagestyle{fancy}
% LE: left even
% RO: right odd
% CE, CO: center even, center odd
% My name for when I print my lecture notes to use for an open book exam.
% \fancyhead[LE,RO]{Gilles Castel}

\fancyhead[R]{\thetitle} % Right odd,  Left even
\fancyhead[L]{\theauthor}          % Right even, Left odd

\fancyfoot[R]{Page \thepage{} of {}\thelastpage}  % Right odd,  Left even
\fancyfoot[L]{\textit{ \nouppercase{\leftmark}}}          % Right even, Left odd
\fancyfoot[C]{}     % Center

\makeatother

%%%%%%%%%%%%%%%%%%%%
% Custom commands	 %
%%%%%%%%%%%%%%%%%%%%
\newcommand{\R}{\mathbb{R}}
\newcommand{\Q}{\mathbb{Q}}
\newcommand{\C}{\mathbb{C}}
\newcommand{\N}{\mathbb{N}}
\newcommand{\Z}{\mathbb{Z}}
\newcommand{\E}{\mathbb{E}}
\newcommand{\V}{\mathbb{V}}
\renewcommand{\P}{\mathbb{P}}

\newcommand{\B}{\mathcal{B}}
\newcommand{\M}{\mathcal{M}}
\renewcommand{\L}{\mathcal{L}}
\newcommand{\F}{\mathcal{F}}
\newcommand{\T}{\mathcal{T}}
\renewcommand{\H}{\mathcal{H}}
\renewcommand{\S}{\mathcal{S}}

\newcommand{\X}{\mathfrak{X}}
\newcommand{\Y}{\mathfrak{Y}}

\newcommand{\powerset}{\mathcal{P}}

\newcommand\nulvec{\bm{0}}
\newcommand\nulmat{\bm{O}}
\newcommand{\what}[1]{\widehat{#1}}
\DeclareMathOperator{\sign}{sign}
\DeclareMathOperator{\sgn}{sign}
\DeclareMathOperator{\re}{Re}
\DeclareMathOperator{\im}{Im}
\DeclareMathOperator{\cov}{Cov}
\DeclareMathOperator{\var}{VaR}
\DeclareMathOperator{\es}{ES}
\newcommand{\1}{\mathds{1}}
\renewcommand{\d}{\text{ d}}
\newcommand{\e}{\mathrm{e}}

% Put x \to \infty below \lim
\let\svlim\lim\def\lim{\svlim\limits}
\let\supup\sup
\let\infup\inf
\let\svsup\sup\def\sup{\svsup\limits}
\let\svinf\inf\def\inf{\svinf\limits}

%Make implies and impliedby shorter
\let\implies\Rightarrow
\let\impliedby\Leftarrow
\let\iff\Leftrightarrow

% switch commands for epsilon
\let\temp\epsilon
\let\epsilon\varepsilon
\let\varepsilon\temp

%Switch commands for phi
% \let\temp\phi
% \let\phi\varphi
% \let\varphi\temp
%

% Define norm and abs to be height adjusted
\DeclarePairedDelimiter\abs{\lvert}{\rvert}%
\DeclarePairedDelimiter\norm{\lVert}{\rVert}%
\DeclarePairedDelimiter\ceil{\lceil}{\rceil}
\DeclarePairedDelimiter\floor{\lfloor}{\rfloor}


% Swap the definition of \abs* and \norm*, so that \abs
% and \norm resizes the size of the brackets, and the 
% starred version does not.
\makeatletter
\let\oldabs\abs
\def\abs{\@ifstar{\oldabs}{\oldabs*}}
%
\let\oldnorm\norm
\def\norm{\@ifstar{\oldnorm}{\oldnorm*}}
\makeatother

\makeatletter
\let\oldfloor\floor
\def\floor{\@ifstar{\oldfloor}{\oldfloor*}}

\makeatletter
\let\oldceil\ceil
\def\ceil{\@ifstar{\oldceil}{\ceil*}}


% Add \contra symbol to denote contradiction
\usepackage{stmaryrd} % for \lightning
\newcommand\contra{\scalebox{1.5}{$\lightning$}}


% Add independent symbol
\newcommand\independent{\protect\mathpalette{\protect\independenT}{\perp}}
\def\independenT#1#2{\mathrel{\rlap{$#1#2$}\mkern3mu{#1#2}}}

\newcommand{\indp}{\independent}


